The acronym \textbf{SAT} comes from the word \textbf{satisfiability}. A SAT problem is a well-known decision problem: given a propositional formula is 
satisfiable or not? \vspace{7pt}

Before moving on, we define formally what is a propositional formula and some properties about it.
\begin{definition}
    A formula in propositional logic is composed by boolean variables and logical operators, which are divided into:
    \begin{itemize}
        \renewcommand{\labelitemi}{-}
        \item $\neg \rightarrow$ negation
        \item $\vee \rightarrow$ disjuction
        \item $\wedge \rightarrow$ conjuction
        \item $\rightarrow \rightarrow$ implication
        \item $\leftrightarrow \rightarrow$ double implication
    \end{itemize} 
    For example, the following formulas:
    $$ (p \vee q) $$
    $$ \neg (p \wedge q) \vee (r \vee t) $$
    are propositional formulas.
\end{definition}
\begin{definition}
    A propositional formula is \textbf{satisfiable} if and only if it is possible to find a truth assignment to the variables that yields the formula true.
\end{definition}
\begin{example}
    e.g. Examples of satisfiability and unsatisfiability. \vspace{7pt}

    For $p = T$ and $q = F$, the following propositional formulas are:

    $$ (p \vee q) \wedge (\neg p \vee \neg q) \wedge (p \vee \neg q) \rightarrow \text{ SAT} $$ 
    $$ (p \vee q) \wedge (\neg p \vee \neg q) \wedge (p \vee \neg q) \wedge (\neg p \vee q) \rightarrow \text{ UNSAT} $$

    We evaluate each time the whole formula. When we prove a formula we said that it is \textbf{SAT} (\textit{satisfiable}) or whether we are not able to 
    satisfy it we said that it is \textbf{UNSAT} (\textit{unsatisfiable}).
\end{example}